\newpage
\begin{section}{Transmitter}

    \begin{subsection}{Method}
    The transmitter is responsible for sending signals from the train, making it possible to locate and track the train's position remotely.
    The transmitter recives power from the track (see Section~\ref{sec:Power supply from track}) and uses the power to generate a electromagnetic frequency.
    The frequency generated is then picked up by recivers and the amplitude of the frequency is used to triangulate the trains possition.

    \end{subsection}
    \begin{subsection}{Implementation}
        The frequency is generated using a ESP32 and a H-bridge hooked up to a coil. 
        The power from the track is rectified using a full bridge rectifier with a capacitor, 
        then sent to a buck converter to steady and lower the voltage for the H-bridge

    \end{subsection}
    \begin{subsection}{Results}

    \end{subsection}
    \begin{subsection}{Discussion}

    \end{subsection}

\end{section}